% !TeX spellcheck = ru_RU
% !TEX root = vkr.tex

\section{Постановка задачи} \label{sec:task}

\subsection{Формальная постановка}

Пусть $\mathcal{X}$~--- пространство пользовательских запросов на естественном языке, поступающих AI-агенту в предметной области поиска объектов недвижимости. Определим конечное множество инструментов (tools):
\[
    \mathcal{T} = \{t_{\mathrm{GEO}},\; t_{\mathrm{TAG}}\},
\]
где $t_{\mathrm{GEO}} \equiv \texttt{ResolveEntities}$ --- инструмент разрешения географических сущностей (город, район, улица, адрес), а $t_{\mathrm{TAG}} \equiv \texttt{SearchTags}$ --- инструмент поиска по характеристикам объекта (площадь, тип ремонта, этажность, количество комнат).

\begin{definition}
Классификатором первого вызова инструмента называется отображение
\[
    f_\theta \colon \mathcal{X} \to \mathcal{T},
\]
реализуемое большой языковой моделью (LLM) с параметрами $\theta$, которое каждому входному запросу $x \in \mathcal{X}$ ставит в соответствие инструмент $t \in \mathcal{T}$, выбираемый моделью в качестве первого вызова инструмента (\texttt{tool\_call}) в генерируемом ответе.
\end{definition}

В качестве базовой линии (baseline) для сравнения используется модель YandexGPT~5.1~Pro ($f_{\theta_0}$), не подвергавшаяся дообучению. Дообученные модели $f_{\theta^*}$ получаются из YandexGPT Lite~--- компактной версии семейства YandexGPT~--- методом LoRA (Low-Rank Adaptation) через Yandex Cloud API.

Введём ошибки классификации в терминах проверки статистических гипотез. Примем за нулевую гипотезу $H_0$: <<запрос $x$ относится к классу TAG>>.

\begin{definition}
\emph{Ошибка I рода}: запрос класса TAG классифицирован как GEO, то есть $f_\theta(x) = t_{\mathrm{GEO}}$ при истинной метке $y(x) = t_{\mathrm{TAG}}$. Такая ошибка приводит к тому, что агент пытается найти несуществующий географический объект вместо применения фильтра по характеристикам.
\end{definition}

\begin{definition}
\emph{Ошибка II рода}: запрос класса GEO классифицирован как TAG, то есть $f_\theta(x) = t_{\mathrm{TAG}}$ при истинной метке $y(x) = t_{\mathrm{GEO}}$. Такая ошибка приводит к потере географической привязки в результатах поиска.
\end{definition}

Дополнительно выделяется третий исход: модель не генерирует вызов инструмента вовсе, что является критическим сбоем в контексте AI-агента. Обозначим долю таких ответов как $\rho_{\varnothing}$.

\subsection{Гипотезы исследования}

\begin{hypothesis}\label{hyp:lora_significance}
(H1a, значимость улучшения.) Дообучение модели YandexGPT~Lite методом LoRA статистически значимо снижает долю ошибок классификации первого вызова инструмента на~бенчмарке \texttt{benchmark-combined} по~сравнению с~базовой моделью YandexGPT~5.1~Pro (accuracy~$= 81{,}8\%$). Проверка~--- односторонний точный критерий Макнемара, $\alpha = 0{,}05$.
\end{hypothesis}

\begin{hypothesis}\label{hyp:lora_threshold}
(H1b, превышение порога.) Нижняя граница 95\%-го bootstrap-доверительного интервала accuracy лучшей дообученной модели на~бенчмарке \texttt{benchmark-combined} превышает~$0{,}90$, что соответствует требованию прикладной системы к~точности классификации.
\end{hypothesis}

\begin{hypothesis}\label{hyp:data_quality}
Качество обучающих данных (сбалансированный микс реальных пользовательских запросов и синтетических примеров) оказывает большее влияние на итоговое качество классификации, чем объём обучающей выборки. В частности, модель, обученная на меньшем, но качественно подготовленном датасете, превосходит модель, обученную на большем, но несбалансированном датасете.
\end{hypothesis}

\subsection{Исследовательские вопросы}

Для систематической проверки сформулированных гипотез определим следующие исследовательские вопросы:

\begin{description}
    \item[\textsc{RQ1}.] \emph{Каково влияние дообучения на метрики классификации?} --- Количественная оценка изменений accuracy, macro~F1, precision и recall по классам $t_{\mathrm{GEO}}$ и $t_{\mathrm{TAG}}$ между базовой моделью $f_{\theta_0}$ и дообученными моделями $f_{\theta^*}$ на бенчмарках \texttt{benchmark-combined} (500 примеров: 326~GEO / 174~TAG) и \texttt{benchmark-real} (250 примеров: 203~GEO / 47~TAG).

    \item[\textsc{RQ2}.] \emph{Как зависит качество классификации от объёма обучающих данных?} --- Анализ зависимости метрик от размера обучающей выборки при фиксированном составе и качестве данных.

    \item[\textsc{RQ3}.] \emph{Как тип обучающих данных влияет на робастность модели?} --- Сравнение моделей, обученных на синтетических данных, реальных пользовательских запросах и их комбинациях, с точки зрения обобщающей способности и устойчивости к изменению распределения входных данных.

    \item[\textsc{RQ4}.] \emph{Статистически значимы ли различия между моделями?} --- Применение критериев Макнемара и bootstrap-доверительных интервалов для попарного сравнения моделей.
\end{description}

\subsection{Критерии успеха}

Критерии успеха определяются на основании baseline-метрик базовой модели YandexGPT~5.1~Pro и целевых значений, обоснованных требованиями прикладной системы. Результаты сведены в Таблице~\ref{tab:success_criteria}.

\begin{table}[H]
\centering
\caption{Критерии успеха: baseline-метрики базовой модели и целевые значения}
\label{tab:success_criteria}
\def\arraystretch{1.2}
\setlength\tabcolsep{0.4em}
\begin{tabular}{l r r l}
    \toprule
    \textbf{Метрика} & \textbf{Baseline} & \textbf{Цель} & \textbf{Бенчмарк} \\
    \midrule
    Accuracy           & $81{,}8\%$ & $> 90\%$  & combined \\
    Macro F1           & $0{,}526$  & $> 0{,}85$ & combined \\
    Recall ($t_{\mathrm{TAG}}$) & $49{,}4\%$ & $> 75\%$  & combined \\
    Без tool\_call ($\rho_{\varnothing}$) & $6{,}8\%$  & $< 3\%$   & combined \\
    \bottomrule
\end{tabular}
\end{table}

Низкий Recall для класса $t_{\mathrm{TAG}}$ и~дисбаланс Macro~$F_1$ (Табл.~\ref{tab:success_criteria}) отражают систематическое предпочтение базовой моделью класса GEO и являются основной мотивацией для дообучения.

\subsection{Ограничения исследования}

Сформулируем ключевые ограничения, в рамках которых проводится данное исследование:
\begin{enumerate}
    \item \textbf{Фиксированная архитектура модели.} Исследование ограничено семейством моделей YandexGPT~5.1 (Lite и Pro). Результаты не обязаны переноситься на модели других архитектур и провайдеров.

    \item \textbf{Бинарная постановка.} Рассматривается только два инструмента из множества инструментов реальной системы. Вопрос масштабирования на $|\mathcal{T}| > 2$ выходит за рамки работы.

    \item \textbf{Дисбаланс классов.} Оба бенчмарка содержат существенное преобладание класса GEO (65\% в \texttt{combined}, 81\% в \texttt{real}), что отражает реальное распределение пользовательских запросов, но ограничивает статистическую мощность выводов по классу TAG.

    \item \textbf{Управляемость LoRA.} Внутренние гиперпараметры дообучения (ранг $r$, $\alpha$, learning rate, расписание, число шагов) управляются платформой Yandex Cloud и недоступны для прямого задания. Это ограничивает возможность выводов о влиянии конкретных гиперпараметров и воспроизводимость результатов на других инфраструктурах.

    \item \textbf{Привязка к версии базовой модели.} Результаты получены на YandexGPT~5.1 в период февраль--март 2026~г. Стабильность результатов при обновлении базовой модели провайдером не гарантируется.

    \item \textbf{Межэкспертное согласие разметки.} Первичная разметка выполнена автором с консенсусным разрешением спорных случаев; формальный коэффициент Cohen's~$\kappa$ не вычислялся ввиду ограниченной доступности независимых разметчиков.
\end{enumerate}
